\clearpage
\appendix
\section{经典微分算子}
\label{chap:appendix-classic-differential-operators}
\label{app:vector-operator-identities}
\renewcommand{\thestatement}{A.1.\arabic{statement}}
\renewcommand{\theHstatement}{appA.sec01.\arabic{statement}}
\setcounter{statement}{0}
\renewcommand{\theequation}{A.\arabic{equation}}
\renewcommand{\theHequation}{appA.\arabic{equation}}
\setcounter{equation}{0}

本附录非常简要地回顾与基本微分算子有关的若干定义和公式, 即梯度, 散度, 旋度和 Laplace 算子. 全书用 $A\cdot B$ 表示 $\mathbb R^d$ 中两个向量的标量积, 并在固定一个定向后, 用 $A\wedge B$ 表示 $\mathbb R^3$ 中两个向量的向量积.

\subsection{标量和向量情形}
\label{sec:app-a-scalar-vector-cases}

\subsubsection{定义}
\label{subsec:app-a-definitions}

\begin{Definition}
\label{def:app-a-basic-vector-operators}
设 $f$ 是标量场, $g$ 是向量场, 二者都足够正则并定义在 $\mathbb R^d$ 的区域 $\Omega$ 上.
\begin{itemize}
 \item $f$ 的梯度是如下向量场: 在 $\mathbb R^d$ 的任意标准正交基中, 它的坐标定义为
 \[
  \nabla f=
  \begin{pmatrix}
   \dfrac{\partial f}{\partial x_1}\\[3pt]
   \vdots\\[3pt]
   \dfrac{\partial f}{\partial x_d}
  \end{pmatrix}.
 \]
 \item $g$ 的散度是定义为如下形式的标量场
 \[
  \operatorname{div}g=\sum_{i=1}^d\frac{\partial g_i}{\partial x_i},
 \]
 其中, 在该空间的一个标准正交基中, $g$ 的坐标记为 $g_1,\ldots,g_d$.
 \item $f$ 的 Laplace 算子是定义为如下形式的标量场
 \[
  \Delta f=\operatorname{div}(\nabla f)
  =\sum_{i=1}^d\frac{\partial^2f}{\partial x_i^2}.
 \]
 \item 当 $d=3$ 时, $g$ 的旋度是记为 $\operatorname{curl}g$ 的向量场, 它在同一标准正交基中的坐标定义为
 \[
  \operatorname{curl}g=
  \begin{pmatrix}
   \dfrac{\partial g_3}{\partial x_2}-\dfrac{\partial g_2}{\partial x_3}\\[5pt]
   \dfrac{\partial g_1}{\partial x_3}-\dfrac{\partial g_3}{\partial x_1}\\[5pt]
   \dfrac{\partial g_2}{\partial x_1}-\dfrac{\partial g_1}{\partial x_2}
  \end{pmatrix}.
 \]
\end{itemize}
此外, 若 $v$ 是 $\mathbb R^d$ 的向量场, 则用 $v\cdot\nabla$ 表示一阶微分算子, 其定义为
\[
 (v\cdot\nabla)f=\sum_{i=1}^d v_i\frac{\partial f}{\partial x_i}.
\]
\end{Definition}

上面定义的概念实际上是内蕴的, 即它们不依赖于基的选择. 因而, 可以在非 Cartesian 坐标中表示这些算子. 回顾 $\mathbb R^3$ 中的柱坐标定义为
\[
 x=r\cos\theta,\qquad y=r\sin\theta,\qquad z=z,
\]
并且该坐标系配有一个旋转的标准正交标架, 记为 $(e_r,e_\theta,e_z)$, 定义为
\[
 e_r=\begin{pmatrix}\cos\theta\\ \sin\theta\\ 0\end{pmatrix},
 \qquad
 e_\theta=\begin{pmatrix}-\sin\theta\\ \cos\theta\\ 0\end{pmatrix},
 \qquad
 e_z=\begin{pmatrix}0\\ 0\\ 1\end{pmatrix}.
\]

\begin{Proposition}
\label{prop:app-a-cylindrical-formulas}
设 $f$ 是标量场, $g$ 是向量场, 二者都足够正则, 并且以柱坐标定义在 $\mathbb R^3$ 的柱形区域 $\Omega$ 上. 则有如下公式
\[
 \nabla f=\frac{\partial f}{\partial r}e_r
 +\frac1r\frac{\partial f}{\partial\theta}e_\theta
 +\frac{\partial f}{\partial z}e_z,
\]
\begin{equation}
 \operatorname{div}g=\frac1r\frac{\partial}{\partial r}(rg_r)
 +\frac1r\frac{\partial g_\theta}{\partial\theta}
 +\frac{\partial g_z}{\partial z},
 \label{eq:app-a-cylindrical-divergence}
\end{equation}
其中 $g=g_re_r+g_\theta e_\theta+g_ze_z$, 并且
\begin{equation}
 \Delta f=\frac1r\frac{\partial}{\partial r}
 \left(r\frac{\partial f}{\partial r}\right)
 +\frac1{r^2}\frac{\partial^2f}{\partial\theta^2}
 +\frac{\partial^2f}{\partial z^2},
 \label{eq:app-a-cylindrical-laplacian}
\end{equation}
以及
\[
 \operatorname{curl}g=
 \left(\frac{\partial g_z}{\partial\theta}-\frac{\partial g_\theta}{\partial z}\right)e_r
 +\left(\frac{\partial g_r}{\partial z}-\frac{\partial g_z}{\partial r}\right)e_\theta
 +\left(\frac1r\frac{\partial(rg_\theta)}{\partial r}
 -\frac1r\frac{\partial g_r}{\partial\theta}\right)e_z.
\]
\end{Proposition}

\subsubsection{常用公式}
\label{subsec:app-a-useful-formulas}

设 $A,B,C$ 是三个向量场, $f$ 是标量场. 先用如下称为向量三重积的代数公式开始这一小组公式:
\begin{equation}
 A\wedge(B\wedge C)=(A\cdot C)B-(A\cdot B)C.
 \label{eq:app-a-vector-triple-product}
\end{equation}
接下来有如下公式, 它们把上面引入的各个微分算子联系起来:
\begin{align}
 \operatorname{div}(fA)&=f(\operatorname{div}A)+(A\cdot\nabla)f,
 \label{eq:app-a-div-scalar-vector}\\
 \operatorname{curl}(fA)&=f(\operatorname{curl}A)+(\nabla f)\wedge A,
 \label{eq:app-a-curl-scalar-vector}\\
 (A\cdot\nabla)A&=\frac12\nabla|A|^2-A\wedge\operatorname{curl}A,
 \label{eq:app-a-convective-identity}\\
 \operatorname{div}(A\wedge B)&=B\cdot\operatorname{curl}A-A\cdot\operatorname{curl}B,
 \label{eq:app-a-div-vector-product}\\
 \operatorname{curl}(A\wedge B)&=(\operatorname{div}B)A-(\operatorname{div}A)B
 +(B\cdot\nabla)A-(A\cdot\nabla)B,
 \label{eq:app-a-curl-vector-product}\\
 \operatorname{curl}(\operatorname{curl}A)&=\nabla\operatorname{div}A-\Delta A.
 \label{eq:app-a-curl-curl}
\end{align}

当然, 这个列表远非完备, 它仅包含流体力学中最有用的公式.
